\documentclass[11pt]{article}
\usepackage{shared/hanzocover}
\usepackage[utf8]{inputenc}
\usepackage[T1]{fontenc}
\usepackage{amsmath,amssymb}
\usepackage{graphicx}
\usepackage{booktabs}
\usepackage{xcolor}
\usepackage{listings}
\input{shared/lstlang}
\usepackage[hidelinks]{hyperref}

\definecolor{codegreen}{rgb}{0,0.6,0}
\definecolor{codegray}{rgb}{0.5,0.5,0.5}
\definecolor{codepurple}{rgb}{0.58,0,0.82}
\definecolor{backcolour}{rgb}{0.96,0.96,0.94}
\definecolor{hzblue}{rgb}{0.12,0.30,0.55}
\definecolor{hzred}{rgb}{0.60,0.16,0.16}

\lstdefinestyle{mystyle}{
    backgroundcolor=\color{backcolour},
    commentstyle=\color{codegreen},
    keywordstyle=\color{codepurple},
    numberstyle=\tiny\color{codegray},
    stringstyle=\color{codegreen},
    basicstyle=\ttfamily\footnotesize,
    breakatwhitespace=false,
    breaklines=true,
    captionpos=b,
    keepspaces=true,
    numbers=none,
    showspaces=false,
    showstringspaces=false,
    showtabs=false,
    tabsize=2,
    frame=single,
    rulecolor=\color{codegray}
}
\lstset{style=mystyle}

\title{The Native Object Store and Its POSIX Projection:\\
       Measured Throughput, Latency, and the Limits of Filesystem Emulation}
\author{Hanzo AI Research\thanks{Corresponding author: research@hanzo.ai} \\
    \textit{Hanzo Industries Inc (Techstars '17)} \\
    Los Angeles, California \\
    \texttt{https://hanzo.ai}}
\date{Revision~1}

\begin{document}

\hanzocoverpage

\twocolumn[
  \begin{@twocolumnfalse}
  \maketitle
  \begin{abstract}
  \noindent
  We report measurements of \texttt{hanzoai/s3}, the native Hanzo object
  store, and of the POSIX filesystem it projects through a
  Container Storage Interface (CSI) driver whose control path speaks the
  native ZAP transport rather than gRPC. All figures are taken from the
  production deployment, in-cluster, against authenticated requests.

  The object store is fast: \textbf{2.26\,ms} per read on a reused
  connection, 265.9\,MiB/s on 16\,MiB objects, and stable latency under
  mixed concurrent load. The filesystem projection sustains 287\,MB/s
  sequential writes and 685\,MB/s sequential reads.

  Its metadata plane does not follow, and we locate the cost precisely. The
  store creates files at \textbf{120/s} (8.36\,ms) when driven directly over
  its own API; the same operation through the FUSE mount costs \(\approx\)
  37\,ms, so \textbf{the projection layer adds \(4.4\times\)}, not the
  store. On a real metadata-bound workload---a Git repository---the mount is
  \(27\times\) to \(278\times\) slower than local disk.

  We therefore do not recommend relocating a Git forge onto this mount as it
  stands, while being explicit that \emph{this is a tuning ceiling we have
  not established}: the FUSE layer's parameters are unexplored, and the
  store's own 8.36\,ms create is itself unexamined. The durable boundary we
  do claim is narrower: \textbf{bulk sequential data projects well into
  POSIX; dependent metadata chains pay a per-operation round trip that must
  be measured before it is designed around.}

  We also document three defects found while bringing this path to
  production---two in connection-pooled RPC and one silently destroying data
  during metadata restore---because each was invisible to the tests that
  would ordinarily be trusted, and the reasons they were invisible are
  reusable.
  \end{abstract}
  \vspace{1em}
  \end{@twocolumnfalse}
]

\section{System Under Test}

\texttt{hanzoai/s3} is a fork of the SeaweedFS lineage that has replaced its
filer RPC surface with ZAP, the Hanzo capability transport: 28 unary and 5
streaming methods answer over \texttt{transport.ListenStream} on the port
historically named the ``gRPC port''. No protobuf framing remains on that
path. The measured build is \texttt{v1.0.4}.

The metadata store is \texttt{luxdb}, a wrapper over
\texttt{github.com/luxfi/database} with a \texttt{zapdb} backend. This is the
sole store compiled into the binary; the SeaweedFS \texttt{leveldb2} store was
removed. That removal is consequential and we return to it in
Section~\ref{sec:migration}.

The POSIX projection is \texttt{hanzoai/s3-csi}, a CSI driver forked so that
its filer-facing half speaks ZAP. A volume is a directory beneath the filer
root, so volumes are \texttt{ReadWriteMany} by construction. The
kubelet-facing half remains CSI gRPC, which is the Kubernetes contract rather
than a design choice.

\subsection{Deployment Shape, Stated Honestly}
\label{sec:durability}

It is worth being exact about durability, because the architecture is often
described more ambitiously than the deployment warrants.

\begin{itemize}
  \item Volume replication is configured \texttt{"000"}: \textbf{one copy of
        every volume}. There is no redundancy at the store layer.
  \item The topology contains \textbf{one} volume server and one master
        member.
  \item The Deployment is \texttt{Recreate}, \texttt{replicas: 1}, bound to a
        single 700\,GiB \texttt{ReadWriteOnce} block volume.
  \item \texttt{luxfi/consensus} is \textbf{not} a dependency of this
        service. No consensus protocol participates in its durability.
\end{itemize}

Durability therefore rests entirely on the underlying cloud block device and
its snapshots. The store is not highly available: a roll is an outage, and a
lost volume is a restore-from-backup event. We state this plainly because a
paper that implied otherwise would be describing a system we do not run.
Introducing replication, or placing the metadata plane under a consensus
protocol, remains future work rather than shipped behaviour.

\section{Method}

Measurements were taken from pods inside the same cluster as the store, to
exclude ingress and WAN effects. Object-store figures use SigV4-signed
requests through the service DNS name; the driver and the store were the
production instances serving live traffic. Filesystem figures use
\texttt{dd} with \texttt{conv=fsync} for writes.

\paragraph{A correction, and the reason it matters.} Our first pass reported
per-operation latencies roughly \(7\times\) too high. It invoked a fresh
client process per request, so each measurement included a DNS resolution and
a TCP handshake. A breakdown of one \(15.2\,\)ms request showed
\(9.0\,\)ms of name lookup, \(2.0\,\)ms of connect, and only
\(4.1\,\)ms of server work. The tell was visible in the data before we
looked for it: a 1\,KiB PUT measured \emph{slower} than a 64\,KiB PUT, which
no server-side cost model explains. Figures below use a single reused
connection and are reported per operation.

The Git comparison runs an identical script against the CSI mount and against
the pod's local disk, so the delta isolates the storage path.

\section{Object Store Results}

\begin{table}[h]
\centering
\begin{tabular}{lrrr}
\toprule
Object size & PUT (ms) & GET (ms) & GET (MiB/s) \\
\midrule
1\,KiB   &  38.93 & 15.67 &   0.1 \\
64\,KiB  &  18.21 & 17.10 &   3.7 \\
1\,MiB   &  36.46 & 20.30 &  49.3 \\
16\,MiB  & 244.80 & 60.17 & 265.9 \\
\bottomrule
\end{tabular}
\caption{Signed request latency, in-cluster, mean of 10.}
\end{table}

The table above retains the per-connection figures, which are what a client
that does not pool connections actually experiences. With a \textbf{reused
connection}, 200 sequential reads average \textbf{2.263\,ms} each
(max 20.3\,ms), so roughly \(13\,\)ms of the naive figure is connection
setup rather than service time. Any client in this cluster should pool.

Under a sustained mixed load---eight rounds of four concurrent
1000-key \texttt{LIST} operations interleaved with
\texttt{PUT}/\texttt{GET}/\texttt{DELETE}---the store completed every
operation with no failures and unary latency flat between 7.5 and
21\,ms. This load shape is precisely what exposed the transport defect
of Section~\ref{sec:hol}, and its stability is the evidence that the defect
is resolved.

\section{Filesystem Projection Results}

\begin{table}[h]
\centering
\begin{tabular}{lrr}
\toprule
Operation & Block size & Rate \\
\midrule
Sequential write & 1\,MiB & 188\,MB/s \\
Sequential write & 4\,MiB & 287\,MB/s \\
Sequential read  & 1\,MiB & 651\,MB/s \\
Sequential read  & 4\,MiB & 685\,MB/s \\
\midrule
File creation & --- & \textbf{27 files/s} \\
\texttt{stat} & --- & \textbf{201 stats/s} \\
\bottomrule
\end{tabular}
\caption{CSI mount, 20\,GiB RWX volume.}
\end{table}

The bulk-data path is strong and would satisfy most streaming workloads. The
namespace path is not: creating a file costs approximately 37\,ms and
\texttt{stat} approximately 5\,ms.

\paragraph{Decomposing the cost, layer by layer.} We measured the same
logical operation---create one metadata entry---at three depths, using a
purpose-built client that speaks ZAP directly to the filer.

\begin{table}[h]
\centering
\begin{tabular}{lrr}
\toprule
Path & Create & Lookup \\
\midrule
ZAP, RAM-backed store   & \textbf{1.41\,ms} & \textbf{1.14\,ms} \\
ZAP, block-backed store & 11.39\,ms & 6.61\,ms \\
Filer HTTP, block store &  8.36\,ms & --- \\
Through the FUSE mount  & \(\approx\)37\,ms & \(\approx\)5\,ms \\
\bottomrule
\end{tabular}
\caption{Identical operation at four depths. Only the backing store differs
between rows one and two: same binary, same transport, same client.}
\end{table}

Three things follow, and they redirect the entire optimisation effort.

\textbf{First, the transport is not the bottleneck.} ZAP's own floor is
1.41\,ms per create and 1.14\,ms per lookup---712 and 879 operations per
second. It also beats the filer's HTTP surface (5.97--11.39\,ms vs
8.36\,ms) on the same storage. The binary protocol is doing its job.

\textbf{Second, the store is disk-sync bound by roughly \(8\times\).}
Moving the data directory from network block storage to \texttt{tmpfs} takes
creates from 11.39\,ms to 1.41\,ms with no other change. Entries are written
one \texttt{Put} at a time; a batch path exists in the store but the
single-entry insert does not use it, so every metadata write pays a separate
durable commit to a network block device. \emph{Lookups} show the same shape
(6.61\,ms vs 1.14\,ms), meaning hot entries are not being served from memory
either.

\textbf{Third, only then does FUSE matter}---and it matters a lot, adding
roughly \(3\times\) on top of the block-backed store. Since each filesystem
operation costs several filer round trips, any reduction in per-round-trip
latency is multiplied by that fan-out. That makes the store-layer fix the
higher-leverage one even though FUSE is the larger single term.

\paragraph{A tuning attempt we are not counting.} We enabled the driver's
file-chunk read cache, which ships disabled. It is inconclusive and we report
it as such: the knob caches \emph{chunks} for reads while this workload is
bound by \emph{metadata}, our two runs did not time identical scopes, and
mount daemons restarted mid-measurement. It was the wrong knob, chosen before
the decomposition above existed.

\section{The Metadata Cost, Made Concrete}

Aggregate rates understate how punishing this is, because real workloads
issue metadata operations in dependent chains rather than in parallel. We
therefore measured Git, which is metadata-bound almost by definition: loose
objects, refs, and index updates are small files created and stat'ed in
sequence.

\begin{table}[h]
\centering
\begin{tabular}{lrrr}
\toprule
Operation & Local & CSI mount & Slowdown \\
\midrule
init + 200 files + commit & 1617\,ms & 43496\,ms & \(27\times\) \\
\texttt{clone}            & 83\,ms   & 23044\,ms & \(278\times\) \\
\texttt{gc --aggressive}  & 162\,ms  & 9699\,ms  & \(60\times\) \\
\texttt{log} + \texttt{status} & 12\,ms & 2857\,ms & \(238\times\) \\
\bottomrule
\end{tabular}
\caption{Identical Git workload, same pod, two storage paths.}
\end{table}

A 278-fold penalty on \texttt{clone} is not a tuning problem. It is
the arithmetic of thousands of dependent round trips at 5--37\,ms
apiece.

\subsection{Consequence for the Git Forge}

The motivating goal of this work was to relocate a 22\,GiB Git forge
(329 bare repositories) off a single \texttt{ReadWriteOnce} volume---the
constraint pinning the control-plane binary to one node---onto shared
storage. The forge's storage layer is an \texttt{osfs} rooted at a data
directory, so the change appeared to require no code at all: mount the
directory elsewhere.

The measurements say do not. Git's heavy paths (clone/fetch serve, push
receive) shell out to the \texttt{git} CLI against an OS path and stream
multi-gigabyte packfiles, which this mount serves well. But every other Git
operation is metadata, and metadata is where the projection collapses.

The honest conclusions are:
\begin{enumerate}
  \item \textbf{Do not place the live Git forge on this mount.} The
        \(238\times\) penalty on \texttt{log}/\texttt{status} alone would be
        user-visible on every request.
  \item \textbf{The object store is the right home for packfiles}, which are
        large, sequential, and immutable---exactly the shape it serves at
        265.9\,MiB/s.
  \item \textbf{Decouple the two problems.} Single-node pinning is a
        \emph{placement} problem; solving it with a filesystem projection
        imports a metadata cost the workload cannot absorb. A Git-native
        object backend, or replication at the block layer, addresses
        placement without paying it.
\end{enumerate}

We regard this as the paper's main result, and note it inverts the plan we
began with. The prior reasoning---that the code was already correct and only
the mount was missing---was true about the code and wrong about the
workload.

\section{Defects Found En Route}

Three defects had to be fixed before any of the above could be measured.
Each is reported because each defeated the check that should have caught it.

\subsection{Promise identifiers scoped to the client, not the connection}

Generated ZAP clients each carried a private \texttt{rpc.Session} numbering
calls from~1, while the connection pool shares one connection per peer
address and constructs a fresh client per call. Two concurrent calls
therefore both claimed promise~1; the second registration evicted the first
from the connection's in-flight table, and the evicted response was discarded
on arrival. The caller blocked for the lifetime of its context---at service
start-up, \texttt{context.Background()}.

The invariant: \textbf{a promise identifier is the connection's
demultiplexing key, so the connection must allocate it}. Stream identifiers
already did.

\subsection{Metadata restore that dropped every file}

The filer's \texttt{meta.load} created directories inline and files on
goroutines, and returned on end-of-input without waiting for them. The caller
closes the connection on return, so in-flight writes died against a torn-down
connection---and the command reported success. A restored tree came back with
its directory structure intact and \emph{no files}.

This is the failure mode one discovers during a disaster recovery. It
survived because the only signal was the absence of data nobody had yet
tried to read.

\subsection{Head-of-line blocking on pooled connections}
\label{sec:hol}

Inbound stream frames were delivered from the connection's read loop with a
blocking send into a 16-slot buffer. A consumer slower than its producer
filled that buffer, stalled the read loop, and with it every other stream and
every unary reply on the connection. Pooling turned a throughput concern into
an outage: one large \texttt{LIST} froze every request in the process.

It presented as an authentication fault, because unauthenticated requests are
rejected before reaching the filer and continued answering in
5\,ms. Diagnosis required a goroutine dump from the stalled process; no
log line existed, because the code was not failing, only waiting.

The fix queues per stream and wakes the reader, so delivery never blocks the
read loop; a bound fails one stream rather than the connection everyone
shares.

\subsection{A configuration that could not say ``I don't know''}

The change that exposed the fsync cost also broke the filesystem mount, and
the mechanism generalises past this codebase.

The mount's metadata cache supplies its own implementation of the store
configuration interface, and that implementation returned its own directory
path for \emph{every} key that did not end in \texttt{backend}. Reading a
newly-added \texttt{syncWrites} option therefore yielded
\texttt{/tmp/<id>/meta}, the store correctly rejected it as neither true nor
false, and the mount died at startup. The store was right; the configuration
was wrong to invent an answer.

The defect is not the one key. A configuration object that cannot express
absence will mis-answer whichever option is added next, and it does so
silently at the call site that is furthest from the mistake. The fix is that
unknown keys return empty, and the regression test asserts the \emph{silence}
rather than any particular key.

Two contributing details are worth naming because both are habits rather than
bugs. The reading side used a variadic parameter, which made a positional
mistake compile; it is now explicit, so a caller that forgets the argument
fails to build. And the interface offers no ``is this set?'' predicate, which
is what pushed a three-state question (absent / true / false) through a
two-state type---the same shape as the downstream bug in
Section~\ref{sec:defects-common}, one layer up.

\subsection{What they have in common}
\label{sec:defects-common}

Two of the four were latent properties of \emph{connection pooling}---a
sharing optimisation that converts an isolated stall into a global one---and
two were three-state questions forced through two-state types, where
``unset'' and ``false'' became indistinguishable. All were invisible to
unauthenticated smoke tests, to fresh-dataset tests, and to green unit suites.
The tests that found them were: an adversarial reproduction with a
deliberately slow consumer; a restore exercised against data written by the
\emph{previous} release; a goroutine dump taken while production was hung; and
an A/B harness built to measure something else entirely.

\section{Store Migration}
\label{sec:migration}

Because \texttt{luxdb} is the only store compiled in, a version upgrade
against an existing \texttt{leveldb2} data directory starts with an
\textbf{empty namespace}---no error, no crash, every bucket invisible. The
two layouts coexist in one directory (\texttt{leveldb2} in numbered
sub-directories, \texttt{zapdb} in top-level files), which is what makes the
migration reversible.

The production migration moved 50{,}006 directories and 68{,}646 files, and
all 45 buckets returned. Reversal is an image change; only writes made in the
new store are lost. The upgrade path is therefore: export, roll, import,
verify with a \emph{signed} operation against \emph{pre-existing} data---the
last qualifier being the one that matters, since unauthenticated and
freshly-created-object checks pass while the store is broken.

\section{Conclusion}

The native object store is fast---2.26\,ms per pooled read---and its POSIX
projection is fast for bulk data and slow for namespaces. But the asymmetry is
\emph{not} simply the inherent cost of emulating a filesystem over a network
metadata service, and we were wrong to frame it that way in an earlier draft.
Measured layer by layer, the store creates at 120/s and the mount at 27/s: most
of the gap is the projection layer, which is implementation subject to tuning
we have not attempted.

What survives as a claim is narrower and, we think, more useful: a dependent
metadata chain pays one round trip per link, and no amount of bulk throughput
compensates for that. But the per-link cost is not a constant of the design.
Decomposed, it is 1.41\,ms of transport, plus \(\approx 10\,\)ms of durable
commit to a network block device, plus a FUSE multiplier on top. Two of those
three terms are addressable without changing the architecture at all:
group-commit the metadata writes (the batch path already exists, unused by the
single-entry insert), and serve hot entries from memory rather than disk.

That is the difference between a workload being impossible and merely
expensive, and it is why the number has to be measured per layer before an
architecture is built on an assumption about it. We began this work believing
the object store was too slow for a filesystem projection. It is not: it
sustains 712 creates per second when it is not waiting on a disk.

Bulk, sequential, immutable data belongs in the object store. Namespaces with
dependent metadata chains---Git foremost---do not, and should not be moved
there merely because their storage layer makes it syntactically easy.

\end{document}
